\section{Обговорення}
\label{sec:discussion}

\subsection{Квазістатична сім'я станів та ініціювання відгалуження}
\label{subsec:discussion-quasistatic}

У розглядуваній постановці V-подібна міжфазна тріщина зсуву є заданою
стаціонарною конфігурацією: довжина кожної її гілки $l$ не змінюється, а
зі зміною параметра зовнішнього навантаження $C$ змінюється рівноважна
довжина локальної зони передруйнування $d_i$. Тому розв'язок задає
однопараметричну квазістатичну сім'ю станів
\begin{equation*}
  C\longrightarrow d_i(C)
  \longrightarrow
  \bigl\{\delta_i(C),W_i(C),J_i(C)\bigr\}.
\end{equation*}
Квазістатичність тут означає послідовність рівноважних конфігурацій зони
за відсутності інерційних ефектів.

Ця сім'я станів безпосередньо придатна для визначення критичного
навантаження ініціювання відгалуженої тріщини. Якщо незалежно задано
критичне розкриття $\delta_{i\mathrm{c}}$ або критичний зональний
енергетичний параметр $J_{i\mathrm{c}}$, то критичний стан визначається
умовою
\begin{equation*}
  \delta_i(C_c)=\delta_{i\mathrm{c}}
  \qquad\text{або}\qquad
  J_i(C_c)=J_{i\mathrm{c}}.
\end{equation*}
Після незалежного визначення відповідної критичної характеристики
матеріалу ці співвідношення безпосередньо задають $C_c$. Подальше
поширення та траєкторія вже сформованої скінченної відгалуженої тріщини
потребують окремої постановки.

\subsection{Конкурентні механізми руйнування}
\label{subsec:discussion-competing}

У стаціонарній V-подібній тріщині можливі два локально різні механізми.
Перший пов'язаний з її кутовою точкою $O$: сингулярне поле породжує
розтягувальну зону передруйнування в одному з матеріалів, а досягнення
критичного $\delta_i$ або $J_i$ визначає ініціювання відгалуження.
Другий механізм локалізований біля зовнішніх вершин V-подібної тріщини.
Для них у роботі~\cite{NazarenkoKipnis2022Equilibrium} визначено
коефіцієнт інтенсивності напружень $K_{\mathrm{II}}$ та сформульовано
умову міжфазного руйнування
\begin{equation*}
  K_{\mathrm{II}}=K_{\mathrm{IIc}}.
\end{equation*}

Отже, відгалуження з кутової точки та міжфазне поширення із зовнішніх
вершин слід розглядати як конкурентні механізми в одній і тій самій
геометричній конфігурації, а не як наперед задані послідовні стадії.
Який із них реалізується при меншому зовнішньому навантаженні,
визначається зіставленням відповідних критичних навантажень: для
зовнішніх вершин --- за міжфазною тріщиностійкістю моди~II, для кутової
точки --- за критичною характеристикою матеріалу, що відповідає
$\delta_i$ або $J_i$.

Маломасштабне припущення $d_i\ll l$ дає змогу розглядати ці локальні
області у провідному наближенні окремо. Поле стаціонарної V-подібної
тріщини використовується як зовнішнє для задачі про зону біля $O$, а
зворотний вплив малої зони на поля біля зовнішніх вершин не
враховується. Для кількісного аналізу конкуренції механізмів за $d_i/l$,
що вже не є малим параметром, потрібна зв'язана крайова задача.

\subsection{Енергетичний і деформаційний критерії}
\label{subsec:discussion-criteria}

Відповідно до співвідношення~\eqref{eq:energy-opening-relation}, для
фіксованої геометрії та пружних характеристик
\begin{equation*}
  J_i=\chi_i\sigma_i\delta_i.
\end{equation*}
Тому $J_i$ і $\delta_i$ не є незалежними характеристиками поточного
стану зони: обидві величини однозначно визначаються тією самою
рівноважною довжиною $d_i$. Звідси деформаційний критерій
$\delta_i=\delta_{i\mathrm{c}}$ і енергетичний критерій
$J_i=J_{i\mathrm{c}}$ вибирають один і той самий критичний стан за
узгодженого співвідношення
\begin{equation*}
  J_{i\mathrm{c}}
  =\chi_i\sigma_i\delta_{i\mathrm{c}}.
\end{equation*}

Ця еквівалентність є властивістю даної однопараметричної моделі та
фіксованої геометрії. Вона не означає універсальної тотожності двох
незалежно визначених критичних характеристик для довільних задач,
оскільки коефіцієнт $\chi_i$ залежить від кутової конфігурації та
пружного контрасту. Практично це означає, що для прогнозування
ініціювання можна користуватися тією критичною характеристикою, яка
надійніше визначається для даного матеріалу.

\subsection{Залишкове кутове поле та локальний енергетичний потік}
\label{subsec:discussion-flux}

У конфігурації із зоною передруйнування поле в кутовій точці $O$
залишається сингулярним. Якщо $H_i$ --- амплітуда провідного кутового
власного поля, то
\begin{equation*}
  \sigma\sim H_i r^{\lambda_i}.
\end{equation*}
Відповідна густина пружної енергії має порядок
\begin{equation*}
  w\sim\frac{H_i^2}{E_i}r^{2\lambda_i},
\end{equation*}
а локальний конфігураційний потік енергії через контур характерного
радіуса $R$ масштабується як
\begin{equation}
\label{eq:local-elastic-flux-scaling}
  J_i^{(\mathrm{el})}(R)
  \sim\frac{H_i^2}{E_i}R^{2\lambda_i+1}.
\end{equation}

Для фізично допустимих режимів, отриманих у числовому аналізі,
$-1/2<\lambda_i<0$. Тому $2\lambda_i+1>0$ і
\begin{equation*}
  \lim_{R\to0}J_i^{(\mathrm{el})}(R)=0.
\end{equation*}
Отже, залишкова необмеженість напружень у точці $O$ не породжує
скінченного ненульового локального пружного внеску до енергетичного
параметра зони. У цьому сенсі кутова сингулярність у конфігурації із
зоною передруйнування слабша за класичну тріщинну сингулярність
$r^{-1/2}$.

Саме тут проходить межа застосовності адитивного представлення,
використаного в роботі~\cite{KaminskyDudyk2026Criteria}. Для
розглянутої там уже існуючої тріщини зовнішнє поле має показник з дійсною частиною $-1/2$, і
скінченний пружний енергетичний потік може входити до загального
енергетичного балансу разом із внеском зони передруйнування. Для
розглядуваної кутової точки за $\lambda_i>-1/2$ локальна межа
$J_i^{(\mathrm{el})}(R)$ прямує до нуля, тому окремий скінченний
доданок $J_i^{(\mathrm{el})}$ до критерію ініціювання не потрібний.
Зональний енергетичний параметр $J_i=\dd W_i/\dd d_i$ тому може бути
використаний без такого адитивного поправочного члена; це не надає йому
статусу незалежно доведеного глобального контурного інтеграла.

\subsection{Область застосовності та перспективи узагальнення}
\label{subsec:discussion-scope}

Отримані співвідношення ґрунтуються на розділенні масштабів
$d_i\ll l\ll L$. Основною областю їх кількісного застосування є
$d_i/l<0.1$. Ділянки числових кривих, де $d_i/l$ перевищує $0.1$ і
поблизу максимумів досягає приблизно $0.18$, корисні для відображення
тенденцій, але вже є екстраполяцією локального наближення. У цих режимах
зростає роль наступних членів зовнішньої асимптотики та взаємного впливу
локальних областей біля кутової точки й зовнішніх вершин.

Локальний характер розв'язку водночас спрощує його перенесення на
ширший клас задач. Для розвинутої симетричної V-подібної міжфазної
тріщини з довільною довжиною гілок множник $CQ_i$, що задає амплітуду
зовнішнього поля для локальної задачі, може бути замінений амплітудою
(узагальненим коефіцієнтом інтенсивності) відповідного кутового
сингулярного поля. Таке узагальнення залишається коректним за умови
маломасштабності зони та збереження тієї самої локальної власної задачі.
Це відкриває можливість використовувати модель для прогнозування
ініціювання відгалуження в куті зламу межі поділу. Частковим випадком
може бути задача розколювання твердого тіла гострим клином, якщо її
локальна крайова постановка зводиться до аналогічної кутової
конфігурації.
